% GENERATED FILE. Edit canonical lessons, then run npm run generate:books.
% canonical-source-hash: fnv1a64:005796aea7c77fdc
% canonical-lessons: IT-C297-arrestare, IT-C297-arresto, IT-C297-manette, IT-C297-catturare, IT-C297-fuggire

\chapter{In manette}
\label{ch:it-in-manette}
\hlchaptermodality{297}

\begin{chapteropening}
\textbf{By the end of this chapter:} \emph{I can report an arrest: arresting, the handcuffs, the capture and running away.}

Complete the last lesson of chapter 297: I can report an arrest: arresting, the handcuffs, the capture and running away.
\end{chapteropening}

\section[arrestare]{arrestare — to arrest}

\label{lesson:IT-C297-arrestare}



\begin{quote}
Before the new one: say the Italian for to question, then the Italian for to confess.
\end{quote}

\subsection*{You'll want to know: arrestare}
\textbf{arrestare} — \textquotedblleft{}to arrest\textquotedblright{}.

\begin{quote}
Hanno arrestato due persone.
\end{quote}

\begin{quote}
\emph{They arrested two people.}
\end{quote}

In formal Italian also ‘to stop’: \emph{arrestare la corsa}.

\begin{cousinweb}
\emph{arrestare} is usually traced to Vulgar Latin \emph{arrestāre} \textquotedblleft{}to make stop\textquotedblright{}, built on Latin \emph{restāre} \textquotedblleft{}to stay behind\textquotedblright{} (Lewis \& Short).
\end{cousinweb}

\begin{grammarlens}[title={What you've built}]
One more word for telling, explaining and describing at B1.
\end{grammarlens}

\subsection*{Your turn}
\begin{itemize}
\raggedright
  \item \emph{Say it:} \emph{arrestare}
  \item \emph{Say it:} \emph{arrestare}, once more
  \item \emph{Say it:} a sentence of your own with \emph{arrestare}
  \item \emph{Recall:} say the Italian for to question, then the Italian for to confess, then say \emph{arrestare} again
\end{itemize}

\begin{tcolorbox}[breakable,colback=teal!4,colframe=teal!35!black,title={Before you move on}]
What does arrestare mean? (\textquotedblleft{}To arrest\textquotedblright{}.) Say it once more.
\end{tcolorbox}
\section[l'arresto]{l'arresto — the arrest; the stop}

\label{lesson:IT-C297-arresto}



\begin{quote}
Before the new one: say the Italian for to confess, then the Italian for to arrest.
\end{quote}

\subsection*{You'll want to know: l'arresto}
\textbf{l'arresto} — \textquotedblleft{}the arrest; the stop\textquotedblright{}.

\begin{quote}
L'arresto è avvenuto alle tre di notte.
\end{quote}

\begin{quote}
\emph{The arrest took place at three in the morning.}
\end{quote}

\emph{Gli arresti domiciliari} — house arrest.

\begin{cousinweb}
\emph{arresto} is built on \emph{arrestare} \textquotedblleft{}to arrest\textquotedblright{}, from Vulgar Latin \emph{arrestāre}, built on \emph{restāre} \textquotedblleft{}to stay behind\textquotedblright{} (Lewis \& Short).
\end{cousinweb}

\begin{grammarlens}[title={What you've built}]
One more word for telling, explaining and describing at B1.
\end{grammarlens}

\subsection*{Your turn}
\begin{itemize}
\raggedright
  \item \emph{Say it:} \emph{l'arresto}
  \item \emph{Say it:} \emph{l'arresto}, once more
  \item \emph{Say it:} a sentence of your own with \emph{l'arresto}
  \item \emph{Recall:} say the Italian for to confess, then the Italian for to arrest, then say \emph{l'arresto} again
\end{itemize}

\begin{tcolorbox}[breakable,colback=teal!4,colframe=teal!35!black,title={Before you move on}]
What does l'arresto mean? (\textquotedblleft{}The arrest; the stop\textquotedblright{}.) Say it once more.
\end{tcolorbox}
\section[le manette]{le manette — the handcuffs}

\label{lesson:IT-C297-manette}



\begin{quote}
Before the new one: say the Italian for to arrest, then the Italian for the arrest.
\end{quote}

\subsection*{You'll want to know: le manette}
\textbf{le manette} — \textquotedblleft{}the handcuffs\textquotedblright{}.

\begin{quote}
Il poliziotto gli ha messo le manette.
\end{quote}

\begin{quote}
\emph{The policeman put the handcuffs on him.}
\end{quote}

Plural; literally ‘little hands’.

\begin{cousinweb}
\emph{manette} is the plural of \emph{manetta}, a little \emph{mano} \textquotedblleft{}hand\textquotedblright{}, from Latin \emph{manus} (Lewis \& Short).
\end{cousinweb}

\begin{grammarlens}[title={What you've built}]
One more word for telling, explaining and describing at B1.
\end{grammarlens}

\subsection*{Your turn}
\begin{itemize}
\raggedright
  \item \emph{Say it:} \emph{le manette}
  \item \emph{Say it:} \emph{le manette}, once more
  \item \emph{Say it:} a sentence of your own with \emph{le manette}
  \item \emph{Recall:} say the Italian for to arrest, then the Italian for the arrest, then say \emph{le manette} again
\end{itemize}

\begin{tcolorbox}[breakable,colback=teal!4,colframe=teal!35!black,title={Before you move on}]
What does le manette mean? (\textquotedblleft{}The handcuffs\textquotedblright{}.) Say it once more.
\end{tcolorbox}
\section[catturare]{catturare — to capture, to catch}

\label{lesson:IT-C297-catturare}



\begin{quote}
Before the new one: say the Italian for the arrest, then the Italian for the handcuffs.
\end{quote}

\subsection*{You'll want to know: catturare}
\textbf{catturare} — \textquotedblleft{}to capture, to catch\textquotedblright{}.

\begin{quote}
Hanno catturato il rapinatore dopo una settimana.
\end{quote}

\begin{quote}
\emph{They caught the robber after a week.}
\end{quote}

Also for animals: \emph{catturare un leone}.

\begin{cousinweb}
\emph{catturare} is built on \emph{cattura} \textquotedblleft{}capture\textquotedblright{}, from Latin \emph{captūra} \textquotedblleft{}a catching\textquotedblright{}, from \emph{capere} \textquotedblleft{}to take\textquotedblright{} (Lewis \& Short).
\end{cousinweb}

\begin{grammarlens}[title={What you've built}]
One more word for telling, explaining and describing at B1.
\end{grammarlens}

\subsection*{Your turn}
\begin{itemize}
\raggedright
  \item \emph{Say it:} \emph{catturare}
  \item \emph{Say it:} \emph{catturare}, once more
  \item \emph{Say it:} a sentence of your own with \emph{catturare}
  \item \emph{Recall:} say the Italian for the arrest, then the Italian for the handcuffs, then say \emph{catturare} again
\end{itemize}

\begin{tcolorbox}[breakable,colback=teal!4,colframe=teal!35!black,title={Before you move on}]
What does catturare mean? (\textquotedblleft{}To capture, to catch\textquotedblright{}.) Say it once more.
\end{tcolorbox}
\section[fuggire]{fuggire — to flee, to run away}

\label{lesson:IT-C297-fuggire}



\begin{quote}
Before the new one: say the Italian for the handcuffs, then the Italian for to capture.
\end{quote}

\subsection*{You'll want to know: fuggire}
\textbf{fuggire} — \textquotedblleft{}to flee, to run away\textquotedblright{}.

\begin{quote}
Il rapinatore è fuggito dalla finestra.
\end{quote}

\begin{quote}
\emph{The robber fled through the window.}
\end{quote}

More formal than \emph{scappare}: \emph{è fuggito}.

\begin{cousinweb}
\emph{fuggire} is inherited from Latin \emph{fugere} \textquotedblleft{}to flee\textquotedblright{} (EtymDB), through a Vulgar Latin \emph{fugīre}.
\end{cousinweb}

\begin{grammarlens}[title={What you've built}]
One more word for telling, explaining and describing at B1.
\end{grammarlens}

\subsection*{Your turn}
\begin{itemize}
\raggedright
  \item \emph{Say it:} \emph{fuggire}
  \item \emph{Say it:} \emph{fuggire}, once more
  \item \emph{Say it:} a sentence of your own with \emph{fuggire}
  \item \emph{Recall:} say the Italian for the handcuffs, then the Italian for to capture, then say \emph{fuggire} again
\end{itemize}

\begin{tcolorbox}[breakable,colback=teal!4,colframe=teal!35!black,title={Before you move on}]
What does fuggire mean? (\textquotedblleft{}To flee, to run away\textquotedblright{}.) Say it once more.
\end{tcolorbox}
